problem:  
**Problem (Noncommutative Rigidity of Hörmander Hamiltonian Flows).**  
Let \( (M,\mathcal H,g_{\mathcal H}) \) be a compact smooth sub‑Riemannian manifold whose horizontal distribution \(\mathcal H\subset TM\) satisfies Hörmander’s bracket‑generating condition with constant rank.  
Let \(H(x,p)=\tfrac12\sum_{i=1}^k \langle p,X_i(x)\rangle^2\) be the sub‑Riemannian Hamiltonian on \(T^*M\) induced by an orthonormal frame \(\{X_1,\dots,X_k\}\), and let  
\[
\Phi_t: T^*M\setminus 0 \longrightarrow T^*M\setminus 0
\]
be its Hamiltonian (horizontal geodesic) flow.

Define
\[
A_{\mathcal H}:=C_0(T^*M\setminus 0)\rtimes_{\Phi}\mathbb R,
\]
the **crossed‑product \(C^*\)-algebra of the Hörmander Hamiltonian flow**, i.e. the \(C^*\)-algebra of the transformation groupoid associated with \(\Phi_t\).

---

**Conjecture (Hörmander Flow Rigidity Conjecture, Hamiltonian Form).**  
Suppose \((M_i,\mathcal H_i,g_{\mathcal H_i})\), \(i=1,2\), are compact Hörmander sub‑Riemannian manifolds.  
Assume that there exists a *‑isomorphism*
\[
\Phi: A_{\mathcal H_1} \;\longrightarrow\; A_{\mathcal H_2}
\]
which induces the canonical homeomorphism of unit spaces \(T^*M_1\setminus 0 \leftrightarrow T^*M_2\setminus 0\) and intertwines the natural \(\mathbb R\)-actions from the flows.

Does it follow that
\[
(M_1,\mathcal H_1,g_{\mathcal H_1}) \;\text{and}\; (M_2,\mathcal H_2,g_{\mathcal H_2})
\]
are isometric (i.e., that there exists a diffeomorphism \(\varphi\) carrying \(\mathcal H_1\) isometrically onto \(\mathcal H_2\))?  

If rigidity fails in full, does the isomorphism class of \(A_{\mathcal H}\) still determine the **osculating nilpotent type** of the filtered Lie algebra generated by \(\mathcal H\)?

---

**Why it is a "good" problem:**  
This corrected formulation is mathematically sound: the action \(\Phi_t\) is a genuine Hamiltonian flow on \(T^*M\), hence defines a well‑posed crossed‑product \(C^*\)-algebra.  
The conjecture asks whether the purely noncommutative object \(A_{\mathcal H}\)—encoding the orbit structure of the horizontal geodesic flow—determines the underlying filtered geometry.  

It lies at the intersection of **sub‑Riemannian geometry**, **microlocal Hamiltonian dynamics**, and **noncommutative topology**.  
A positive result would yield a new *reconstruction theorem* for hypoelliptic geometry, showing that Hörmander’s analytic data can be recovered from the groupoid algebra of its Hamiltonian dynamics.  
A negative result would uncover “noncommutatively isospectral” sub‑Riemannian manifolds, revealing hidden analytic degeneracies among distinct Hörmander structures.  

The statement is conceptually simple—rigidity of a \(C^*\)-algebra associated with a flow—but resolving it demands unifying techniques from analysis, geometry, and operator algebras, making it a truly deep and elegant **good research problem**.